\section{Límites de alcance}

El sistema está pensado para que una sesión termine bien o termine; no para
que sobreviva a la pérdida de un participante. Conviene decir con precisión
qué significa eso, porque es la diferencia entre una limitación conocida y una
promesa incumplida.

\emph{No hay reconexión.} Un nodo que pierde la conexión no vuelve. No existe
un procedimiento para que se reincorpore con su estado, ni un identificador
que le permita reclamar lo que tenía. El conjunto de participantes se fija al
abrir la sesión y solo puede reducirse.

\emph{No hay recuperación de escrituras no publicadas.} Lo que un nodo escribió
en su marco local y todavía no liberó existe únicamente en su memoria. Si el
proceso muere con un permiso de escritura tomado, ese trabajo se pierde entero
y no hay forma de reconstruirlo: el coordinador nunca lo vio.

\emph{No hay persistencia.} Los datos confirmados viven en la memoria del
coordinador. Cuando la sesión termina, el contenido de las páginas desaparece.
No hay volcado a disco, ni bitácora de escrituras, ni manera de retomar una
sesión anterior.

\emph{No hay tolerancia a particiones.} Una red que se divide no produce dos
mitades que sigan trabajando: produce nodos que dejan de alcanzar al
coordinador y que, por lo tanto, dejan de poder adquirir nada. No hay quórum,
ni elección, ni reconciliación posterior.

Lo que sí hay es un corte limpio, y se distingue según cómo ocurra la pérdida.

Si el coordinador recibe una señal de terminación, avisa: le encola un error
de sesión abortada a cada nodo vivo, vacía las salidas con un tope, y recién
entonces cierra. Cada trabajador recibe un código que dice qué pasó, en vez de
descubrirlo cuando una lectura falla.

Si el coordinador muere sin ejecutar una instrucción más, no hay aviso posible
y tampoco hace falta inventarlo: el núcleo cierra sus sockets, el cierre llega
a cada trabajador y todos lo ven en el acto. Ésa es toda la garantía que se
puede dar en ese caso, y alcanza para que nadie quede esperando.

Del lado del trabajador, un error de sesión es terminal. La biblioteca rechaza
localmente todo lo que puede detectar antes de enviarlo —página fuera de
rango, modo inválido, adquisición anidada—, así que un error que llega desde
el coordinador a un trabajador correcto solo puede venir de un defecto. Seguir
operando sobre un estado que no se entiende es peor que terminar: lo que se
corrompe es la memoria compartida, y el síntoma aparecería lejos de la causa.

Hay además tres límites de tamaño que no son fallas sino decisiones fijadas en
tiempo de compilación: la página mide 4096 bytes, la sesión admite a lo sumo
cuatro nodos, y la cantidad de páginas se fija al arrancar y no cambia. El
primero es el que más se nota en los resultados, porque es la unidad de
coherencia y por lo tanto el tamaño del falso compartimiento. Comparar otro
tamaño de página era una de las ampliaciones posibles del enunciado y no se
hizo; el código la admite —el tamaño es una constante de compilación y no hay
ningún valor literal que dependa de ella— pero las pruebas sí dan por sentado
el tamaño actual, de modo que una construcción con otro valor sirve para medir
y no para correr la suite.
